\documentclass[11pt,a4paper]{article}

\usepackage[T1]{fontenc}
\usepackage[utf8]{inputenc}
\usepackage{lmodern}
\usepackage{amsmath,amssymb,amsthm}
\usepackage{mathtools}
\usepackage{array}
\usepackage{booktabs}
\usepackage{geometry}
\usepackage{microtype}
\usepackage[hidelinks]{hyperref}

\hypersetup{
  pdftitle={A Mobius-binomial C/F bridge and a claim-disciplined theorem-transfer atlas for Apery-like sources},
  pdfauthor={Oleksiy Babanskyy}
}

\geometry{margin=2.5cm}
\pagestyle{plain}

\newtheorem{theorem}{Theorem}[section]
\newtheorem{corollary}[theorem]{Corollary}
\newtheorem{lemma}[theorem]{Lemma}
\newtheorem{remark}[theorem]{Remark}

\DeclareMathOperator{\CT}{CT}
\DeclareMathOperator{\diag}{diag}
\DeclareMathOperator{\ord}{ord}

\title{A M\"obius-binomial C/F bridge and a claim-disciplined theorem-transfer atlas for Ap{\'e}ry-like sources}

\author{Oleksiy Babanskyy}

\date{}

\begin{document}

\maketitle

\begin{abstract}
We isolate a direct M\"obius-binomial bridge between Zagier's C and F
Ap{\'e}ry-like sequences:
\[
  F(x)=(1-9x)^{-1} C\!\left(\frac{-x}{1-9x}\right).
\]
Equivalently, the bridge is an involutive binomial coefficient transform
between the two catalog-normalized coefficient sequences and preserves their
ordinary Hankel determinants; it also transports the differential equation of
C to that of F and is cross-checked by the known Franel transform route.  We
then use this bridge as a primitive for a theorem-transfer atlas.  The atlas
distinguishes closed internal results,
theorem imports by reference, conditional transfers, finite exact replays,
source-attached conjectures, proof obligations, blocked rows, bounded
negatives, watch metadata, and closed metadata.  In particular, we do not
derive Sun's Conjecture 3.5 from the C/F bridge, do not promote finite
replays or coefficient matches to theorems, and do not turn Dwork or
Rowland--Yassawi metadata into unconditional theorem claims.  Sun's raw-C
conjecture is addressed in a separate Franel--Domb bridge manuscript
\cite{babanskyy_2026_sun35_bridge}.  The result here is a claim-disciplined
package for using a proved C/F bridge while keeping adjacent Ap{\'e}ry-like
source layers separate.
\end{abstract}

\section{Introduction}
\label{sec:introduction}

Ap{\'e}ry-like sequences sit at the intersection of several structures:
recurrences, modular parametrizations, binomial transforms, constant-term and
Laurent-polynomial models, supercongruences, and automatic congruences.  These
structures often appear in clusters, but they do not transfer automatically
from one sequence to another.  A coefficient match, a shared recurrence shape,
or a common modular level may suggest a relation; none of these is by itself a
theorem-transfer mechanism.

This point of view is close in spirit to the tabular literature around
Calabi--Yau and Ap{\'e}ry-like differential equations, where large lists of
operators, transformations, and periods are useful only when their
normalizations and equivalences are explicit
\cite{almkvist_van_enckevort_van_straten_zudilin_2005_tables,
almkvist_van_straten_zudilin_2011_clausen}.  Our scope is much narrower: a
claim-disciplined atlas centered on one proved C/F bridge.

This paper is organized around one case where the transfer mechanism is
explicit.  We prove and certify a direct M\"obius-binomial bridge between the
Zagier C and F sequences:
\[
  F(x)=(1-9x)^{-1} C\!\left(\frac{-x}{1-9x}\right).
\]
The identity gives an involutive coefficient transform, an ODE pullback, and a
Hankel determinant invariance, together with a constant-term complement
operation.  It is compatible with the known Franel network around these
sequences.

The second contribution is methodological.  We use the C/F bridge as a
primitive in a theorem-transfer atlas for nearby Ap{\'e}ry-like source layers.
The atlas records what can be transferred, what can only be imported by
reference, what is merely metadata, and what remains a proof obligation.  The
point is not to claim that all adjacent congruence frameworks become theorems
of this paper.  Rather, we separate proof, imported theorem, finite replay,
conjectural source attachment, blocked metadata, watch rows, and bounded
negative evidence.

The bridge is the mathematical core; the atlas records how far its transport
consequences extend.  Sections~\ref{sec:sequences}--\ref{sec:normalization}
give the C/F bridge and its certificate.  Sections
\ref{sec:contracts}--\ref{sec:package} explain how the bridge is used to build
a disciplined theorem-transfer atlas.  Appendices contain the theta-gauge
certificate, the Atkin--Lehner normalization audit, the constant-term
complement guardrail, selected registry rows, open problems, and source-status
notes.

\section{Zagier C and F}
\label{sec:sequences}

Let $C_n$ and $F_n$ denote Zagier's C and F Ap{\'e}ry-like sequences in the
catalog convention used throughout this paper
\cite{zagier_2009_aperylike,oeis_a002893,oeis_a093388}.  Write
\[
  C(x)=\sum_{n\geq0}C_nx^n,\qquad
  F(x)=\sum_{n\geq0}F_nx^n.
\]
The C sequence is given by
\[
  C_n=\sum_{k=0}^n\binom{n}{k}^2\binom{2k}{k}.
\]
It satisfies
\[
  (n+1)^2C_{n+1}
  =(10n^2+10n+3)C_n-9n^2C_{n-1}.
\]
The F sequence is normalized by $F_0=1$, $F_1=6$, and satisfies
\[
  (n+1)^2F_{n+1}
  =(17n^2+17n+6)F_n-72n^2F_{n-1}.
\]
The atlas below always treats these catalog-normalized objects as the named
source and target rows.

\section{The C/F M\"obius-binomial bridge}
\label{sec:bridge}

\subsection{Modular parameter and theta-gauge normalization}
\label{subsec:theta-normalization}

We use the modular parametrizations recorded by Beukers in the common
$\Gamma_0(6)$ setting \cite{beukers_2025_supercongruences_modular_forms}.  In
the C row, let $t_C(q)$ denote the source modular parameter and put
\[
  \Theta_C(q)=C(t_C(q)).
\]
For the F row, Beukers' source parameter is twist-guarded.  We write it as
$t_F^{\rm source}(q)$ and use the catalog parameter
\[
  t_F(q)=t_F^{\rm cat}(q):=-t_F^{\rm source}(q).
\]
The F theta series used below is therefore
\[
  \Theta_F(q)=F(t_F(q))
             =F(-t_F^{\rm source}(q))
             =\Theta_F^{\rm source}(q).
\]
All $q$-expansions, cusp-order checks, and Sturm checks in this paper are in
this common normalization.  The two F parameters are kept distinct:
\[
  t_F^{\rm source}=\frac{t_C}{1-9t_C},\qquad
  t_F^{\rm cat}=-\frac{t_C}{1-9t_C}.
\]
Consequently
\[
  1+9t_F^{\rm source}=1-9t_F^{\rm cat}=\frac{1}{1-9t_C}.
\]
In particular, the two identities used in the bridge proof are identities of
meromorphic modular functions on $X_0(6)$ in this normalization.

\begin{theorem}[Zagier C/F M\"obius-binomial bridge]
\label{thm:cf-bridge}
Under the catalog convention used here and the verified F twist transport,
\[
  F(x)=(1-9x)^{-1} C\!\left(\frac{-x}{1-9x}\right).
\]
\end{theorem}

\begin{proof}
With the normalization of Section~\ref{subsec:theta-normalization}, the
theta-gauge identities are
\[
  \Theta_F/\Theta_C = 1-9t_C,\qquad
  \Theta_C/\Theta_F = 1+9t_F^{\rm source}=1-9t_F^{\rm cat}.
\]
The second form is the one compatible with the catalog parameter.  With
\[
  u=t_C(q),\qquad x=t_F^{\mathrm{cat}}(q)=\frac{-u}{1-9u},
\]
one has
\[
  \frac{-x}{1-9x}=u,\qquad (1-9x)^{-1}=1-9u.
\]
Thus
\[
  F(x)=\Theta_F(q)=(1-9u)\Theta_C(q)
      =(1-9x)^{-1}C\!\left(\frac{-x}{1-9x}\right).
\]
The theta-gauge identities are verified by forward and inverse valence
certificates, with a cleared-denominator Sturm backup; see
Appendix~\ref{app:certificate}.  The same normalization also has an
Atkin--Lehner \(W_2\) parameter and theta audit recorded in
Appendix~\ref{app:al-w2}; this audit is explanatory and is not used as a
substitute for the M\"obius-binomial bridge proof.
\end{proof}

\begin{corollary}[Coefficient transform]
\label{cor:coeff}
Define
\[
  T(a)_m=\sum_{n=0}^m\binom{m}{n}(-1)^n9^{m-n}a_n.
\]
Then $T(C)=F$ and $T(F)=C$.
\end{corollary}

\begin{proof}
Taking coefficients in Theorem~\ref{thm:cf-bridge} gives $T(C)=F$.  Let
$g(x)=-x/(1-9x)$.  Then $g(g(x))=x$ and
\[
  (1-9x)^{-1}(1-9g(x))^{-1}=1.
\]
Applying the same transform a second time therefore gives $T(F)=C$.
\end{proof}

\begin{corollary}[Hankel determinants]
\label{cor:hankel}
For a sequence $A=(A_n)_{n\geq0}$ put
\[
  \Delta_N(A)=\det(A_{i+j})_{0\leq i,j<N}.
\]
Then, for every $N\geq1$,
\[
  \Delta_N(C)=\Delta_N(F).
\]
\end{corollary}

\begin{proof}
This is the standard Hankel invariance of binomial transforms
\cite{spivey_steil_2006_k_binomial}, but the short argument is included to
fix the signs.  Set $G_m=(-1)^mF_m$.  By Corollary~\ref{cor:coeff},
\[
  G_m=\sum_{n=0}^m\binom{m}{n}(-9)^{m-n}C_n.
\]
For fixed $N$, let $B_N$ be the lower triangular matrix
\[
  (B_N)_{i,j}=
  \begin{cases}
    \binom{i}{j}(-9)^{i-j},& j\leq i,\\
    0,& j>i,
  \end{cases}
  \qquad 0\leq i,j<N.
\]
The matrix $B_N$ has determinant $1$, and Vandermonde convolution gives
\[
  (G_{i+j})_{0\leq i,j<N}
  =
  B_N\,(C_{i+j})_{0\leq i,j<N}\,B_N^{t}.
\]
Hence $\Delta_N(G)=\Delta_N(C)$.  Finally,
\[
  (F_{i+j})_{0\leq i,j<N}
  =
  D_N\,(G_{i+j})_{0\leq i,j<N}\,D_N,
  \qquad D_N=\operatorname{diag}((-1)^0,\ldots,(-1)^{N-1}),
\]
so $\Delta_N(F)=\Delta_N(G)$.
\end{proof}

\begin{corollary}[ODE pullback]
\label{cor:ode}
Let
\[
  L_C=x(x-1)(9x-1)D_x^2+(27x^2-20x+1)D_x+3(3x-1)
\]
be the differential operator associated to the C recurrence, and let
\[
  L_F=x(72x^2-17x+1)D_x^2+(216x^2-34x+1)D_x+6(12x-1)
\]
be the corresponding F operator.  Under
\[
  u=-\frac{x}{1-9x},\qquad F(x)=(1-9x)^{-1}C(u),
\]
symbolic substitution gives
\[
  L_C[C](u)=(9x-1)L_F[F](x).
\]
Thus the C equation pulls back to the F equation up to the rational factor
$9x-1$, which is nonzero as a rational function.
\end{corollary}

\begin{remark}[Franel route]
\label{rem:franel}
Let $R_n=\sum_{k=0}^n\binom{n}{k}^3$.  The known Franel transforms are
\[
  C(x)=(1-x)^{-1}R\!\left(\frac{x}{1-x}\right),\qquad
  F(x)=(1-8x)^{-1}R\!\left(\frac{-x}{1-8x}\right),
\]
with provenance through Barrucand, Callan, Verrill, and OEIS
\cite{barrucand_1975_franel_c,callan_2008_franel_c,verrill_2010_congruences_modular_forms}.
Composing these formulas gives the same C/F bridge.
\end{remark}

\section{Certificate and normalization}
\label{sec:normalization}

The proof depends on four normalization checks.  First, C and F are placed in a
common $\Gamma_0(6)$ normalization.  Second, the source F parameter is
transported to the catalog-normalized F twist.  Third, the forward and inverse
theta-gauge identities are certified by cusp-order and valence arguments.
Fourth, a cleared-denominator Sturm check gives an independent backup.  The
main text uses only the resulting identities; the certificate details are
recorded in Appendix~\ref{app:certificate}.

\section{Additional functorial consequences}
\label{sec:contracts}

The C/F bridge is used functorially in a small number of controlled ways.
The closed internal consequences are the coefficient involution, the Hankel
determinant invariance, the ODE pullback, and the binomial-complement
constant-term lemma.  Together with the imported C-side Laurent model, that
lemma gives the $9-\Lambda_{C,1}$ complement witness.  The source-layer
consequences below explain how those facts are recorded in the atlas.  They
do not promote adjacent Rowland--Yassawi, BTY, or Dwork comparison metadata
into unconditional theorem claims.

\subsection{Coefficient transform and Liu normalization}
\label{subsec:coeff-transport}

The coefficient transform in Corollary~\ref{cor:coeff} is the atlas primitive.
It allows a theorem stated for F to be rewritten as a theorem for $T(C)$.  It
does not automatically produce a theorem for raw $C_n$.

In Liu's normalization
\[
  v_m^{(A)}(u)=\sum_{n=0}^m\binom{m}{n}(-u)^{m-n}A_n
\]
\cite{liu_2025_binomial_transforms}, the same bridge reads
$v_m^{(C)}(9)=(-1)^mF_m$ and, by involution,
$v_m^{(F)}(9)=(-1)^mC_m$.  This is a normalization map into Liu's
binomial-transform layer, not a blanket import of every adjacent congruence.

For example, Sun's source rows include congruences for a sequence $Q_n$ with
$Q_n=(-1)^nF_n$ in the relevant convention
\cite{sun_2022_supercongruences}.  The bridge rewrites these as statements
about $(-1)^nT(C)_n$.  Sun's raw-C conjectural congruence is not a consequence
of this C/F bridge; it is addressed separately by a Franel--Domb bridge through
intermediate Franel, Domb, and Rodriguez--Villegas sums
\cite{babanskyy_2026_sun35_bridge}.

\subsection{Franel network}

Remark~\ref{rem:franel} explains why the bridge is not an isolated accident:
C and F are both connected to the Franel sequence by known
M\"obius-binomial transforms.  The theorem-transfer atlas records this as
known source metadata.  It does not treat every inverse or coefficient replay
as a new theorem edge.

The source-layer order is important: the Franel sequence is the common
upstream object, while the C/F bridge is the induced relation after the two
classical transforms are composed.  This is why the paper can use Franel as
a provenance check without making every Franel-side theorem a C/F theorem.

More concretely, both generating functions are obtained from the same cubic
Franel series $R(x)=\sum R_nx^n$ by two different fractional-linear changes of
variable.  The C/F bridge is therefore the composition of two classical
Franel transforms, not a relation inferred only from matching initial
coefficients.  This Franel viewpoint is useful in the atlas because it
separates a known structural network from later theorem-transfer claims.

\subsection{Constant-term complement}

The coefficient transform has a direct constant-term interpretation.

\begin{lemma}[Binomial-complement constant-term transform]
\label{lem:ct-complement}
Let $\Lambda$ be a Laurent polynomial and let
\[
  A_n=\CT(\Lambda^n).
\]
For a scalar $a$, define
\[
  B_n=\sum_{k=0}^n\binom{n}{k}(-1)^ka^{n-k}A_k.
\]
Then
\[
  B_n=\CT((a-\Lambda)^n).
\]
\end{lemma}

\begin{proof}
By the binomial theorem,
\[
  (a-\Lambda)^n=\sum_{k=0}^n\binom{n}{k}a^{n-k}(-\Lambda)^k.
\]
Taking constant terms gives
\[
  \CT((a-\Lambda)^n)
  =\sum_{k=0}^n\binom{n}{k}(-1)^ka^{n-k}\CT(\Lambda^k)=B_n.
\]
\end{proof}

For
\[
  \Lambda_{C,1}:=\Lambda_C=\frac{(x+y+1)(xy+x+y)}{xy}
\]
Gorodetsky records this Laurent model for the C row
\cite{gorodetsky_2023_sporadic_representations}.  Equivalently, if
$P=(x+y+1)(xy+x+y)$, then
\[
  \CT(\Lambda_{C,1}^n)=[x^ny^n]P^n,
\]
and the cited model identifies this coefficient with the catalog-normalized
$C_n$.  Lemma~\ref{lem:ct-complement} gives
\[
  F_n=\CT((9-\Lambda_{C,1})^n).
\]
This is a constant-term operation.  It is not a Laurent mutation claim and it
is not a Newton-equivalence claim between $9-\Lambda_{C,1}$ and other F-side
models such as Gorodetsky's replacement model $Q_F$
\cite{gorodetsky_2023_sporadic_representations}.

\subsection{Rational diagonals and small automata}

The constant-term models also admit rational diagonal contracts, which support
small automatic-congruence metadata for C/F and related rows.  This paper
records those automata as metadata.  It does not claim a full
Rowland--Yassawi theorem-by-reference application, because that would require
the source-level bad-prime and singularity filters in addition to an explicit
Cartier-state construction \cite{rowland_yassawi_2015_automatic_congruences}.

\section{Claim classes and atlas discipline}
\label{sec:claim-classes}

The atlas uses explicit claim classes.  Table~\ref{tab:claim-classes}
summarizes the ones used in this paper.

\begin{table}[htbp]
\centering
\small
\begin{tabular}{>{\ttfamily}p{0.42\linewidth}p{0.48\linewidth}}
\hline
closed\_internal\_result & Proved or certified inside the C/F package.\\
theorem\_import\_by\_reference & External theorem used with its source label; not renamed as an internal theorem.\\
conditional\_theorem\_transfer & Transfer valid after an explicit bridge or source assumption.\\
closed\_metadata & Verified auxiliary record, such as a normalization map, fixed-modulus automaton, source-label reconciliation, or known literature identity; not used as a new theorem unless separately promoted through a proof route.\\
explicit\_automaton\_metadata & Explicit finite-state or prime-power automaton data recorded as metadata; in the V2 package this is treated as a closed-metadata subtype, not as a theorem-transfer class.\\
finite\_exact\_replay & Exact finite check used as evidence only.\\
source\_attached\_conjecture & Conjectural source row attached to the atlas without proof.\\
proof\_obligation & A precise missing lemma, not a vague direction.\\
blocked & A known missing contract prevents theorem use.\\
watch & Interesting metadata without a promotion route.\\
bounded\_negative & A bounded search failed a stated gate.\\
\hline
\end{tabular}
\caption{Claim classes used in the theorem-transfer atlas.}
\label{tab:claim-classes}
\end{table}

The guardrails are simple: finite replay is evidence, not proof; coefficient
matching is not theorem transfer; and external theorems remain external
theorems.

\section{The theorem-transfer registry}
\label{sec:registry}

The registry is an operational table of theorem imports and bridge-generated
transfers.  The complete table is part of the data package; Appendix
\ref{app:registry} records the schema and selected rows.  The core examples
are:
\begin{enumerate}
\item the C/F bridge, a closed internal result;
\item the coefficient involution, a closed internal result;
\item the ODE pullback, a closed internal result;
\item Sun's F-side source rows transferred only to $T(C)$;
\item Sun's raw-C target, retained in the historical V1/V2 atlas as a
  source-attached conjecture and proof obligation, but now marked in the
  paper-facing readout as addressed separately by
  \cite{babanskyy_2026_sun35_bridge};
\item CT/Newton guardrails, which block mutation or Newton-equivalence claims
  without an explicit operation.
\end{enumerate}

\section{Source-layer examples}
\label{sec:source-layers}

\subsection{Sun layer}

Sun supplies source congruences for F-side objects and related binomial sums
\cite{sun_2022_supercongruences}.  The C/F bridge transfers F-side statements
to $T(C)$.  It does not prove raw-C variants from C/F data.  The raw-C
Conjecture 3.5 is a separate Franel--Domb theorem
\cite{babanskyy_2026_sun35_bridge}, and the atlas keeps that proof route
separate from the C/F theorem-transfer rows.

\subsection{Dwork layer}

Mellit--Vlasenko give Dwork-type congruences for constant terms of powers of
Laurent polynomials \cite{mellit_vlasenko_2015_dwork}.  The complement model
$9-\Lambda_{C,1}$ gives an additional F-side constant-term witness.  We record it
as Dwork metadata by reference.  We do not claim a Newton equivalence or a
Laurent mutation between this complement model and the F-side model $Q_F$.

Thus the atlas keeps three objects separate: the C-side Laurent model, the
bridge-generated complement $9-\Lambda_{C,1}$, and independent F-side Laurent
models such as $Q_F$.  The first two give a closed C/F complement
witness; comparing the complement to other F-side models remains open.

\subsection{Automatic congruence layer}

Rowland--Yassawi provide the source framework for automatic congruences of
rational diagonals \cite{rowland_yassawi_2015_automatic_congruences}.  The
atlas records explicit small automata for Cooper, Franel, and C/F small
moduli, including C/F mod $2$, mod $3$, mod $4$, and mod $9$ metadata.  These
are explicit automaton or prime-power metadata; they are not a full
Rowland--Yassawi theorem import.

\subsection{Chan--Cooper companion Lucas lane}

The BTY/Chan--Cooper audit separates the raw weight-one problem from a
weight-two Chan--Cooper companion layer.  The raw weight-one
$\chi(w_d)$ values remain uncomputed.  However, the five companion rows
$V6A,V6B,V6C,V8,V9$ admit direct Chan--Cooper binomial formulas
\cite{chan_cooper_2012_rational_analogues}, and their companion coefficient
sequences satisfy Lucas congruences modulo every prime by a direct
Lucas--Kummer argument.  This result is not a C/F bridge transfer and is not
a BTY Theorem~1.1 import.

Let
\[
  U_i(n)=\binom{2n}{n}s_i(n).
\]
The source formulas are
\[
  s_{6A}(n)=\sum_{j=0}^n(-8)^{n-j}\binom{n}{j}
          \sum_{\ell=0}^j\binom{j}{\ell}^3,
\]
\[
  s_{6B}(n)=\sum_{j=0}^n\binom{n}{j}^2\binom{2j}{j},
  \qquad
  s_{6C}(n)=\sum_{j=0}^n\binom{n}{j}^3,
\]
\[
  s_8(n)=\sum_j4^{n-2j}\binom{n}{2j}\binom{2j}{j}^2,
\]
and
\[
  s_9(n)=\sum_j(-3)^{n-3j}
  \binom{n}{j}\binom{n-j}{j}\binom{n-2j}{j}.
\]
The sums are taken over the natural ranges where the displayed binomial
coefficients are defined.

\begin{theorem}[Chan--Cooper companion Lucas rows]
\label{thm:chan-cooper-companion-lucas}
The five sequences
\[
  U_{6A},U_{6B},U_{6C},U_8,U_9
\]
satisfy Lucas congruences modulo every prime.  The negative-$q$ level-8
sequence
\[
  U_8^-(n)=(-1)^nU_8(n)
\]
also satisfies Lucas congruences modulo every prime.
\end{theorem}

\begin{proof}
We use the standard Lucas theorem for binomial and multinomial coefficients,
together with the Kummer carry criterion.

First, the sequence $a^n$ satisfies Lucas congruences modulo every prime.  If
$p\nmid a$, this follows from Fermat's theorem; if $p\mid a$, the sequence is
congruent to $1,0,0,\ldots$ modulo $p$.  Products of Lucas sequences are
Lucas, and the central binomial sequence $\binom{2n}{n}$ is Lucas, with
digit-doubling carries producing zero factors.

The Franel cubic sum
\[
  R(n)=\sum_j\binom{n}{j}^3
\]
is Lucas by applying Lucas theorem termwise.  This proves the claim for
$s_{6C}$.  The same argument, with the central binomial factor
$\binom{2j}{j}$, proves the claim for
\[
  s_{6B}(n)=\sum_j\binom{n}{j}^2\binom{2j}{j}.
\]
For $s_{6A}$, write
\[
  s_{6A}(n)=\sum_j\binom{n}{j}R(j)(-8)^{n-j}.
\]
This is a binomial convolution of Lucas sequences, hence is Lucas.

For $s_8$, when $p$ is odd, Lucas and Kummer factor the summand
\[
  4^{n-2j}\binom{n}{2j}\binom{2j}{j}^2
\]
digitwise; carry cases vanish modulo $p$.  For $p=2$, one has $s_8(0)=1$
and $s_8(n)\equiv0\pmod2$ for $n>0$, since the $j=0$ term contains $4^n$
and $\binom{2j}{j}$ is even for $j>0$.

For $s_9$, when $p\ne3$, the multinomial form
\[
  (-3)^{n-3j}\frac{n!}{j!j!j!(n-3j)!}
\]
factors digitwise in carry-free cases, and carry cases vanish modulo $p$.
For $p=3$, if $n>3j$, the factor $(-3)^{n-3j}$ is divisible by $3$.  If
$n=3j$ and $j>0$, then $(3j)!/(j!^3)$ is divisible by $3$, since adding
$j+j+j$ in base $3$ creates a carry.  Thus $s_9(0)=1$ and
$s_9(n)\equiv0\pmod3$ for $n>0$.

Multiplication by $\binom{2n}{n}$ preserves Lucas congruences, so the five
$U_i(n)$ are Lucas.  Finally, $(-1)^n$ is Lucas for odd primes and is
invisible modulo $2$, so $U_8^-(n)=(-1)^nU_8(n)$ is Lucas as well.
\end{proof}

This theorem is independent of the C/F bridge.  It is also independent of the
BTY theorem-transfer mechanism.  The level-6 rows have a separate local
BTY/Atkin--Lehner audit, and the $V8,V9$ normalizer geometry remains useful
modular context, but neither is used in the proof above.  No raw weight-one
$\chi(w_d)$ values are computed here.

Other source layers, including Liu's binomial-transform congruence profile,
Straub's Gessel--Lucas theorem, and Beukers--Tsai--Ye's modular-form Lucas
congruence source
\cite{liu_2025_binomial_transforms,straub_2024_gessel_lucas,beukers_tsai_ye_2025_lucas},
are summarized in Appendix~\ref{app:sources}.  They are kept brief here
because their role is source attachment and claim discipline, not the main
C/F theorem.

\section{Open-problem atlas}
\label{sec:open-problems}

The open-problem atlas records blockers rather than vague directions.  A
selection is shown in Table~\ref{tab:open-problems}; a fuller schema appears
in Appendix~\ref{app:open-problems}.

\begin{table}[htbp]
\centering
\small
\begin{tabular}{p{0.25\linewidth}p{0.43\linewidth}p{0.22\linewidth}}
\hline
Branch & Blocking lemma or missing contract & Status\\
\hline
BTY weight-one rows & row-specific $\chi(w_d)$ values & open\\
Straub derivative transfer & low-carry identity for the corrected connection & proof obligation\\
Dwork model comparison & controlled equivalence between complement model and $Q_F$ & open\\
Rowland--Yassawi source application & bad-prime/singularity filter and source-style Cartier states & inactive\\
Edge-watch scan & no unresolved V1 watch rows remain & closed metadata\\
Sun raw-C to RV-108 & Franel--Domb raw-C/RV bridge outside the C/F atlas & companion manuscript\\
A2b global glue & Domb branch-multiplier identity outside the C/F atlas & companion manuscript\\
\hline
\end{tabular}
\caption{Selected open-problem and blocker rows.}
\label{tab:open-problems}
\end{table}

\section{Edge scan and watch-triage discipline}
\label{sec:edge-scan}

A bounded C/F-like M\"obius-binomial scan checked $17684$ transforms of the
form
\[
  B(x)=(1-ax)^{-r}A\!\left(\frac{sx}{1-ax}\right).
\]
The scan used exact integer arithmetic on $12$ coefficients, the $16$ catalog
families recorded in the companion data, $r\in\{1,2,3\}$,
$-200\le a\le200$, and nonzero integers $s$ with $|s|\le20$.
It recovered known or canonical relations and promoted no new theorem edge.
The eight previous watch hits were subsequently resolved as known metadata:
six are alias/canonical identity hits for the Ap{\'e}ry zeta(3) and
Almkvist--Zudilin gamma catalog rows, and two are inverse Franel-transform
hits.  This is a negative discipline result: coefficient matches are not
automatically promoted.

\section{Data package and reproducibility}
\label{sec:package}

The data package contains registry rows, claim classes, source labels, proof
obligations, open-problem statuses, edge-scan triage records, manifests, and
readouts.  The paper-facing purpose is reproducibility of claim status: a
reader should be able to tell which rows are proved internally, which are
imported by reference, which are finite replays, and which are blocked.
The package is based on the V2 manifest, with the V1 registry retained as
immutable provenance for the original C/F theorem-transfer rows.
Appendix~\ref{app:sources} lists source identifiers and package fields.  The
curated companion package is available at
\[
  \text{\url{https://github.com/aconsciousfractal/C-F_Bridge_and_Atlas}},
\]
with the paper source in \texttt{paper/}, JSON artifacts in
\texttt{data/relation\_atlas/}, replay scripts in \texttt{scripts/}, and the
executable-scope note in \texttt{docs/atlas/EXECUTABLE\_REPLAY\_SCOPE.md}.

\section{Conclusion}
\label{sec:conclusion}

We do not solve every problem attached to the C/F bridge.  Instead, we give a
normalized proof of the bridge and a reproducible model for using it without
overclaiming.  The bridge transports generating functions, coefficients, ODEs,
Hankel determinants, and a constant-term complement model.  The atlas records
adjacent theorem sources, metadata layers, finite checks, conjectures,
blockers, and bounded negatives with distinct claim classes.  This separation
is the main lesson:
bridge identities are powerful precisely when their scope is explicit.

\appendix

\section{Theta-gauge certificate details}
\label{app:certificate}

This appendix records the certificate data behind the theta-gauge identities
used in Theorem~\ref{thm:cf-bridge}.

By the normalization contracts in Section~\ref{subsec:theta-normalization},
the functions below are meromorphic modular functions on $X_0(6)$ and have no
interior poles.  Their only possible poles are cuspidal; the cusp-order tables
bound those poles.  The certificate uses the following public data contract.
\[
\begin{array}{c|c|c|c|c}
\text{object} & \text{expression} & \text{weight} & \text{character} & \text{audit source}\\
\hline
t_C & \text{Beukers }X_0(6)\text{ parameter} & 0 & 1 & \text{cusp orders and }q\text{-replay}\\
t_F^{\rm source} & t_C/(1-9t_C) & 0 & 1 & \text{twist-source contract}\\
t_F^{\rm cat} & -t_F^{\rm source} & 0 & 1 & \text{catalog twist contract}\\
\Theta_F/\Theta_C & 1-9t_C & 0 & 1 & \text{forward valence certificate}\\
\Theta_C/\Theta_F & 1-9t_F^{\rm cat} & 0 & 1 & \text{inverse valence certificate}
\end{array}
\]
The checked $q$-expansion orders are recorded here as certificate data.  The
public companion scripts reproduce the finite coefficient and atlas checks;
larger q-expansion/Sturm certificate automation is outside the executable
scope declared in the companion repository.

\subsection{Forward valence certificate}

The forward identity is
\[
  \Theta_F/\Theta_C = 1-9t_C.
\]
Define
\[
  H(q)=\Theta_F(q)/\Theta_C(q)-\bigl(1-9t_C(q)\bigr).
\]
The cusp-order audit on $\Gamma_0(6)$ gives
\[
\begin{array}{c|rrrr}
\text{object} & 0 & 1/2 & 1/3 & \infty\\
\hline
t_C & 0 & -1 & 0 & 1\\
\Theta_F/\Theta_C & 1 & -1 & 0 & 0 .
\end{array}
\]
Thus the only possible pole of $H$ is at the cusp $1/2$, of order at most
$1$.  The $q$-expansion replay gives $\ord_\infty(H)\geq201$.  Since this
exceeds the possible pole order, the valence criterion forces $H=0$.

\subsection{Inverse valence certificate}

The inverse identity, written in the catalog parameter, is
\[
  \Theta_C/\Theta_F = 1-9t_F^{\rm cat}.
\]
Define
\[
  K(q)=\Theta_C(q)/\Theta_F(q)-\bigl(1-9t_F^{\rm cat}(q)\bigr).
\]
The inverse cusp-order audit gives
\[
\begin{array}{c|rrrr}
\text{object} & 0 & 1/2 & 1/3 & \infty\\
\hline
t_F^{\rm cat} & -1 & 0 & 0 & 1\\
\Theta_C/\Theta_F & -1 & 1 & 0 & 0 .
\end{array}
\]
Thus the only possible pole of $K$ is at the cusp $0$, of order at most $1$.
The $q$-expansion replay gives $\ord_\infty(K)\geq201$, and the valence
criterion gives $K=0$.

\subsection{Sturm clear-denominator backup}

The same two identities are also checked after clearing the possible cuspidal
poles.  Use the eta quotient
\[
  M(q)=\frac{\eta(q)^3\eta(2q)^3}{\eta(3q)\eta(6q)}.
\]
By the Ligozat criterion, the exponent vector of $M$ satisfies the two
congruences modulo $24$ and the total weight is $2$.  The cusp orders of $M$
on $\Gamma_0(6)$ are
\[
\begin{array}{c|rrrr}
\text{cusp} & 0 & 1/2 & 1/3 & \infty\\
\hline
\ord(M) & 1 & 1 & 0 & 0 .
\end{array}
\]
Thus $M$ is holomorphic at every cusp and vanishes at the two cusps where a
pole can occur in the preceding valence argument.  The deduction is explicit:
\begin{itemize}
\item $H$ has no possible pole except at $1/2$, of order at most $1$, while
  $\ord_{1/2}(M)=1$; hence $MH$ is holomorphic at $1/2$.
\item At the other cusps $H$ is holomorphic and $M$ is holomorphic, so $MH$
  is holomorphic there as well.
\item $K$ has no possible pole except at $0$, of order at most $1$, while
  $\ord_{0}(M)=1$; hence $MK$ is holomorphic at $0$.
\item At the other cusps $K$ is holomorphic and $M$ is holomorphic, so $MK$
  is holomorphic there as well.
\end{itemize}
Since $H$ and $K$ have weight $0$ and trivial character in the chosen
normalization, both $MH$ and $MK$ are holomorphic weight-$2$ forms on
$\Gamma_0(6)$ with the character of the eta quotient $M$.  The Sturm
criterion used here is the standard one for the corresponding character space;
their checked $q$-expansions lie in the rational coefficient field used by the
certificate, so the ordinary coefficient-vanishing test applies to the cleared
forms.
Since
\[
  [\mathrm{SL}_2(\mathbb Z):\Gamma_0(6)]=12,\qquad
  \left\lfloor\frac{2\cdot12}{12}\right\rfloor=2,
\]
the Sturm bound is $2$.  The replay gives vanishing at infinity far beyond
this bound, hence both cleared forms vanish identically.

\subsection{F twist transport dependency}

The F-side parameter in source notation is twist-guarded.  If
$F(x)=\sum_{n\ge0}F_nx^n$ denotes the catalog-normalized generating function,
the transport contract is
\[
  F\bigl(-t_F^{\rm source}(q)\bigr)=\Theta_F^{\rm source}(q),
\]
or equivalently
\[
  t_F^{\rm cat}(q)=-t_F^{\rm source}(q).
\]
This is why Theorem~\ref{thm:cf-bridge} is stated in catalog-normalized F
notation and why the inverse theta-gauge identity is written as
$\Theta_C/\Theta_F=1-9t_F^{\rm cat}$.

\subsection{Atkin--Lehner W2 parameter and theta audit}
\label{app:al-w2}

The C/F normalization also admits a precise Atkin--Lehner description.  Put
\[
  h=t_F^{\rm source},\qquad t_C=\frac{h}{1+9h},
  \qquad t_F^{\rm cat}=-h,
\]
and let
\[
  W_2=\begin{pmatrix}2&-1\\6&-2\end{pmatrix}.
\]
For weight-zero functions, \(f|W\) means precomposition with \(W\).  For a
weight-one form we use the convention
\[
  (f|_1W)(\tau)=\det(W)^{1/2}(c\tau+d)^{-1}f(W\tau),
\]
with the positive real branch \(\det(W_2)^{1/2}=\sqrt{2}\).

\begin{lemma}[Atkin--Lehner \(W_2\) audit]
\label{lem:al-w2-audit}
In the above normalization,
\[
  t_C|W_2=1+8t_F^{\rm source}=1-8t_F^{\rm cat},
\]
and the theta eta quotients satisfy
\[
  \Theta_C|_1W_2=(2\sqrt{2})^{-1}\Theta_F,
  \qquad
  \Theta_F|_1W_2=-2\sqrt{2}\,\Theta_C.
\]
\end{lemma}

\begin{proof}
The parameter identity follows by substituting the normalized action
\[
  W_2(h)=-\frac{h+1/8}{9(h+1/9)}
\]
into \(t_C=h/(1+9h)\).  This gives \(t_C|W_2=1+8h\).

For the theta statement, write \(\eta_d(\tau)=\eta(d\tau)\).  The eta quotients
are
\[
  \Theta_C=\frac{\eta_2^6\eta_3}{\eta_1^3\eta_6^2},
  \qquad
  \Theta_F=\frac{\eta_1^6\eta_6}{\eta_2^3\eta_3^2}.
\]
The four eta arguments \(dW_2\tau\), for \(d\in\{1,2,3,6\}\), factor
projectively through \(\mathrm{SL}_2(\mathbb Z)\) matrices as follows:
\[
\begin{array}{c|c|c|c}
d & \gamma_d & \text{factorization} & \text{target eta index}\\
\hline
1 & \begin{pmatrix}1&-1\\3&-2\end{pmatrix}
  & \gamma_1\,\diag(2,1) & 2\\
2 & \begin{pmatrix}2&-1\\3&-1\end{pmatrix}
  & \gamma_2\,\diag(1,1) & 1\\
3 & \begin{pmatrix}1&-3\\1&-2\end{pmatrix}
  & \gamma_3\,\diag(6,1) & 6\\
6 & \begin{pmatrix}2&-3\\1&-1\end{pmatrix}
  & \gamma_6\,\diag(3,1) & 3 .
\end{array}
\]
Using the standard Dedekind eta multiplier formula
\[
  \eta(\gamma\tau)=\exp(\pi i R(\gamma))(c\tau+d)^{1/2}\eta(\tau),
  \qquad
  R(\gamma)=\frac{a+d}{12c}-s(d,c)-\frac14,
\]
the exact exponents are
\[
\begin{array}{c|cccc}
d&1&2&3&6\\
\hline
R(\gamma_d)&-1/3&-1/6&-1/3&-1/6 .
\end{array}
\]
For \(\Theta_C\), with eta exponent vector \((-3,6,1,-2)\), the total
root-of-unity exponent is
\[
  -3(-1/3)+6(-1/6)+1(-1/3)-2(-1/6)=0.
\]
For \(\Theta_F\), with exponent vector \((6,-3,-2,1)\), it is
\[
  6(-1/3)-3(-1/6)-2(-1/3)+1(-1/6)=-1.
\]
The remaining \(\tau\)-dependent factors cancel against the weight-one slash
prefactor.  The residual powers of \(2\) are \(2^{-3/2}\) for \(\Theta_C\) and
\(-2^{3/2}\) for \(\Theta_F\), giving the two displayed identities.
\end{proof}

This audit is a normalization statement, not a replacement for the bridge.
Indeed, in the coordinate \(T=t_C\),
\[
  T|W_2=\frac{1-T}{1-9T},
  \qquad
  g(T)=-\frac{T}{1-9T},
\]
so the catalog C/F bridge map is not literally the Atkin--Lehner action.
Rather,
\[
  g(T)=\frac{1-(T|W_2)}{8}.
\]
Thus the F catalog parameter is an affine normalization of the \(W_2\)-translate
of the C parameter.  This observation also does not compute the row-specific
BTY multipliers \(\chi(w_d)\), which remain outside the closed claims of this
paper.

\section{Binomial-complement constant-term lemma}
\label{app:ct-complement}

Lemma~\ref{lem:ct-complement} is included in the main text because it is the
cleanest bridge between coefficient transforms and constant-term models.  The
important guardrail is repeated here: the operation
\[
  \Lambda\mapsto a-\Lambda
\]
preserves the fact that one has a constant-term sequence, but it is not a
Laurent mutation claim and does not imply Newton equivalence between two
chosen Laurent models.

\section{Registry schema and selected rows}
\label{app:registry}

Each registry row uses the following schema:
\[
\begin{array}{ll}
\texttt{statement\_id} & \text{row identifier}\\
\texttt{source\_paper} & \text{external or internal source}\\
\texttt{source\_label} & \text{the theorem, conjecture, or artifact label}\\
\texttt{source\_object} & \text{source sequence or model}\\
\texttt{target\_object} & \text{target sequence or model}\\
\texttt{rewrite} & \text{transport or normalization formula}\\
\texttt{cf\_role} & \text{how the C/F bridge enters}\\
\texttt{claim\_level} & \text{claim class}\\
\texttt{finite\_replay} & \text{finite evidence, if any}\\
\texttt{literature\_status} & \text{source status}\\
\texttt{proof\_obligation} & \text{missing lemma, if any}\\
\texttt{status} & \text{current row status}\\
\texttt{next\_action} & \text{next audit/proof step}.
\end{array}
\]

A representative public excerpt is:
\begin{center}
\scriptsize
\begin{tabular}{@{}p{0.18\linewidth}p{0.12\linewidth}p{0.12\linewidth}p{0.22\linewidth}p{0.22\linewidth}@{}}
\hline
\texttt{statement\_id} & source & target & rewrite/status & claim level\\
\hline
\texttt{cf\_bridge} & C & F & $F=\mathcal{T}(C)$ & closed internal result\\
\texttt{cf\_involution} & C,F & F,C & $\mathcal{T}^2=1$ & closed internal result\\
\texttt{cf\_hankel} & C,F & C,F & $\Delta_N(C)=\Delta_N(F)$ & closed internal result\\
\texttt{cf\_ode\_pullback} & $L_C$ & $L_F$ & rational pullback & closed internal result\\
\texttt{sun\_f\_to\_tc} & Sun F-side & $T(C)$ & signed transfer only & conditional transfer\\
\texttt{sun\_raw\_c} & raw C & RV-108 & no C/F transfer & separate proof obligation\\
\texttt{dwork\_ct} & $\Lambda_{C,1}$ & $9-\Lambda_{C,1}$ & CT complement & closed metadata\\
\texttt{ry\_automata} & small moduli & automata & fixed-modulus data & closed metadata\\
\shortstack[l]{\texttt{chan\_cooper}\\\texttt{five\_companion}\\\texttt{lucas}} & Chan--Cooper & $U_i$ rows & direct Lucas proof & closed internal result\\
\hline
\end{tabular}
\end{center}
The full machine-readable table is the JSON file
\begin{center}
\texttt{data/relation\_atlas/}\\
\texttt{cf\_theorem\_transfer\_registry\_v1.json}.
\end{center}
The excerpt is included only to make the atlas semantics visible in the paper.

\section{Open problems and blocking lemmas}
\label{app:open-problems}

The atlas records open work by the mathematical object that is missing.  For
the Sun branch, the C/F theorem transports F-side statements to $T(C)$, but it
does not itself supply the raw-C/RV-108 bridge.  In the historical V1/V2 data
package this appears as a deferred proof obligation.  After that atlas freeze,
the missing raw-C/RV bridge and its A2b Domb component are recorded in the
separate Franel--Domb bridge manuscript on Sun's Conjecture 3.5
\cite{babanskyy_2026_sun35_bridge}.  We keep the registry row as provenance
for the C/F atlas rather than reclassifying it as a C/F consequence.

The modular-form and constant-term branches have different blockers.  The
Beukers--Tsai--Ye weight-one rows require row-specific multiplier data
$\chi(w_d)$ before they can be used as theorem imports.  The Dwork branch has
a proved C/F constant-term complement operation, but the stronger comparison
between the complement model and the Gorodetsky F-side model $Q_F$ remains an
equivalence-or-obstruction problem.  The Rowland--Yassawi branch has explicit
small automata as metadata, but a full source-theorem application would still
require a bad-prime/singularity filter and source-style Cartier states.

Finally, the Straub derivative-transfer branch is not closed by the formal
C/F coefficient involution alone.  The corrected connection term leaves a
finite low-carry identity as the remaining proof obligation.  The edge-watch
scan no longer contributes an open problem: its eight watch rows have been
resolved as known metadata.

\section{Source and bibliography status}
\label{app:sources}

Zagier and OEIS fix the catalog-normalized C and F rows used throughout the
paper.  Beukers supplies the modular parametrizations, the common
$\Gamma_0(6)$ setting, and the F twist caveat used in the theta-gauge
certificate.  Barrucand, Callan, and Verrill supply the classical Franel
transform context that explains the C/F bridge as part of a known network
rather than as an isolated coefficient coincidence.

The remaining sources enter as theorem imports or metadata layers.  Sun
supplies the source congruence layer and the original Conjecture 3.5 target
\cite{sun_2022_supercongruences}; the latter is addressed in a separate
Franel--Domb bridge manuscript rather than by the C/F bridge
\cite{babanskyy_2026_sun35_bridge}.  Liu supplies binomial-transform
congruence metadata with convention guards \cite{liu_2025_binomial_transforms}.
In the Liu audit, thirteen rows match the normalized catalog directly, while
the $\eta$ and $s18$ rows are stored with an explicit convention map rather
than as failures.  Chan--Cooper supplies the companion binomial formulas for
the five rows in Theorem~\ref{thm:chan-cooper-companion-lucas}; the Lucas proof
there is internal and does not import a BTY theorem
\cite{chan_cooper_2012_rational_analogues}.  Straub's
Gessel--Lucas theorem is imported by reference where its hypotheses are used,
but the separate C/F derivative-transfer problem remains guarded by its
low-carry proof obligation \cite{straub_2024_gessel_lucas}.

The Beukers--Tsai--Ye theorem is used only when row hypotheses are complete
\cite{beukers_tsai_ye_2025_lucas}.  Mellit--Vlasenko and Gorodetsky provide
the Dwork and constant-term context, and Rowland--Yassawi provides the source
contract for automatic congruences
\cite{mellit_vlasenko_2015_dwork,gorodetsky_2023_sporadic_representations,rowland_yassawi_2015_automatic_congruences}.
The broader tabular background of Calabi--Yau and Ap{\'e}ry-like operators is
represented by the Almkvist--van Enckevort--van Straten--Zudilin tables and
the subsequent transformation literature
\cite{almkvist_van_enckevort_van_straten_zudilin_2005_tables,
almkvist_van_straten_zudilin_2011_clausen}.  Before submission, the remaining
bibliographic work is journal-specific style and final access-date policy for
web catalog entries.

\bibliographystyle{alpha}
\bibliography{cf_bridge_relation_atlas_refs}

\end{document}
